% TOP_PROBLEM: {"id": 30006776, "title": "Voevodsky's Smash-Nilpotence Conjecture", "category_id": 5, "category": {"id": 5, "name": "algebraic_geometry", "display_name": "Algebraic Geometry", "description": "Geometric objects defined by polynomial equations.", "slug": "algebraic-geometry", "order_index": 5, "created_at": "2026-07-31T15:26:25.671Z"}, "status": "open"}
% FIRST_POSTED: 2026-09-25
% ENTRY_KIND: statement_only
% !TeX program = lualatex
% Definition and sources only; this file is not an LLM solution attempt.
\documentclass[11pt]{article}
\usepackage[margin=25mm]{geometry}
\usepackage{fontspec}
\setmainfont{Latin Modern Roman}
\usepackage{amsmath,amssymb,mathtools}
\usepackage{unicode-math}
\setmathfont{Latin Modern Math}
\usepackage{xurl,hyperref}
\hypersetup{colorlinks=true,urlcolor=blue}
\setcounter{secnumdepth}{0}
\setlength{\parindent}{0pt}
\setlength{\parskip}{5pt}
\setlength{\emergencystretch}{3em}
\begin{document}
\section*{\texorpdfstring{Voevodsky's Smash-Nilpotence Conjecture}{Voevodsky's Smash-Nilpotence Conjecture}}\label{30006776}
\subsection{Definitions and mathematical statement}
For a smooth projective $d$-dimensional variety over $k$, a rational cycle class $\alpha\in\mathrm{CH}^r(X)_{{\mathbb{Q}}}$ is numerically trivial when
$\deg(\alpha\cdot\beta)=0$ for every complementary class $\beta\in\mathrm{CH}^{d-r}(X)_{{\mathbb{Q}}}$. It is smash-nilpotent when
\[
 \alpha^{\boxtimes N}=0\quad\text{in}\quad\mathrm{CH}^{Nr}(X^N)_{{\mathbb{Q}}}
 \quad\text{for some }N\ge1.
\]
Here $\boxtimes$ is the external product over $k$, and the exponent may depend on $\alpha$. The conjecture says numerical triviality implies smash-nilpotence for every field and every such cycle. Powers under intersection on $X$, or under composition of correspondences, are different operations. Rational coefficients and the absence of a restriction on the characteristic of $k$ are essential scope choices.
\par\smallskip\textit{Formulation and concept references:} \href{https://www.math.ias.edu/vladimir/sites/math.ias.edu.vladimir/files/Nilpotence_theorem_published.pdf}{[S1]}.

\subsection{Short English statement}
Does every rational algebraic cycle invisible to all intersection-number tests become zero after taking sufficiently many external copies of itself?
\subsection{Sources}
\begin{itemize}
\item[S1]Vladimir Voevodsky, A Nilpotence Theorem for Cycles Algebraically Equivalent to Zero, IMRN1995 no.4,187-198; published IAS snapshot,13pages.
\url{https://www.math.ias.edu/vladimir/sites/math.ias.edu.vladimir/files/Nilpotence_theorem_published.pdf}.
\textit{Formulation source.}
\end{itemize}
\end{document}
